\section{The Walk}
\label{sec:ch07-the-walk}

\index{resolvability!is a second pass|(}%
Everything the composer has refused so far it refused while merging the
documents: two claims on one field, decided by reading both declarations. There
is a second pass after that one, and it asks whether the schema the documents
agree on can actually be served.

\index{Ratings subgraph!a third document}%
Here is a third document to ask it with, written by hand like the last one and
exported by nothing. It is the Ratings subgraph
chapter~\ref{ch:federation-model} put on the map, reduced to the two fields
that matter, given a root field of its own so that it is a subgraph rather than
a fragment, and carrying one deliberate mistake:

\begin{graphqlsdl}
schema
  @link(
    url: "https://specs.apollo.dev/federation/v2.6"
    import: ["@key", "FieldSet"]
  ) {
  query: Query
}

type Query {
  ratingCount: Int!
}

type Session @key(fields: "id", resolvable: false) {
  id: ID!
  averageScore: Float
}

scalar FieldSet
\end{graphqlsdl}

\index{resolvable@\code{resolvable: false}!says do not come in here}%
Ratings owns \code{averageScore} and contributes it to a \code{Session} it does
not own, which is the arrangement chapter~\ref{ch:federation-model} said this
example exists to show. The mistake is the argument on the key, and it is a
mistake here rather than in general. \code{resolvable: false} tells the
composer that this subgraph declares the key without being able to answer for
the type, so nothing should try to enter through it. That is the right thing
to say about a type you only point at, and the wrong thing to say about a type
you contribute fields to, because those fields can then be reached from
nowhere. Ratings also exports no \code{node} field, so there is no other door.

\index{composition!merging succeeds and serving would not}%
Merge those two documents and nothing collides. \code{Session} gains a field,
which is exactly what a contributing subgraph is for, and every field in the
result has one owner. The merged schema is fine. It is also unservable, and the
second pass is what says so:

\begin{minted}[fontsize=\footnotesize,breakanywhere]{text}
The field "averageScore" is unresolvable at the following path:
 query {
  node {
   ... on Session {
    averageScore <--
   }
  }
 }
This is because:
 - The root type field "Query.node" is defined in the following subgraph: "sessions".
 - The field "Session.averageScore" is defined in the following subgraph: "ratings".
 - The entity ancestor "Session" in subgraph "sessions" has no accessible target entities (resolvable @key directives)
in the subgraphs where "Session.averageScore" is defined.
 - The type "Session" is not a descendant of any other entity ancestors that can provide a shared route to access
"averageScore".
\end{minted}

\index{resolvability!the error prints the path}%
Read what that error is: a path, printed as a query, with an arrow at the step
the composer could not take. The tool walked from a root field into the schema, kept a record of where
it had got to and which subgraph it was in, and stopped at the first field it
could not reach from there. Then it printed the route it had taken to get
stuck.

\index{figure!the walk}%
Figure~\ref{fig:ch07-the-walk} is that route with the boundary drawn in.

\begin{figure}
  \centering
  % The composer's resolvability walk from Query.node to Session.averageScore,
% as wgc 0.129.9 actually reports it: two steps answered inside subgraph
% sessions, then a third that needs subgraph ratings and has no resolvable
% key to get there on.
%
% Same idiom as figures/tikz/ch01-one-field.tex and
% figures/tikz/ch06-two-routes.tex: a timeline read left to right, ticks
% mark stages, event nodes sit above the axis, and a brace under the
% stage that matters carries a one-line label. This figure has one lane,
% not two, because there is one walk here, not two routes compared
% against each other; the two-lane layout of the earlier figures does not
% apply to a single path that simply stops.
%
% Two things are added to the idiom, both carrying meaning by shape and
% by dashed-versus-solid rather than by colour, per house style:
%
% 1. The axis is solid from tick 1 to tick 2 (the walk stays inside
%    subgraph sessions) and dashed from tick 2 to tick 3 (the step it
%    cannot take). The arrowhead lands on tick 3 and the axis stops
%    there rather than running past it, echoing the composer's own
%    printed diagnostic, which marks the unreachable field with an arrow
%    drawn as <--; this figure does not reproduce that literal arrow but
%    borrows its idea, an arrow that stops at the field it cannot reach.
% 2. Tick 3 is drawn as a cross, not a vertical stroke, marking the point
%    the walk cannot pass, in the manner of a buffer stop.
%
% The brace sits under the dashed segment, at the crossing that does not
% exist, and names the reason: ratings declares the key with
% resolvable: false, so the key is there to be matched against and not to
% be entered by. That argument is the difference between a seam and a
% dead end, and it is why the figure labels the key rather than the type.
%
% No quotation marks anywhere in this file, including this comment block.
% SPEC decision 30 puts the book in strict Ascii mode and the prose gate
% reads a literal quote character as a quotation mark on sight, while the
% escaped form is LaTeX's diaeresis accent and fails to compile inside a
% node. Identifiers are written bare.
%
% Bare tikzpicture (house style): the figure environment, caption and
% label live at the call site in the section file.
\begin{tikzpicture}[
  x=1mm, y=1mm,
  every node/.append style={font=\sffamily\scriptsize},
  axis/.style={draw, line width=0.4pt,
    -{Straight Barb[length=1.4mm, width=1.6mm]}},
  trunk/.style={draw, line width=0.4pt},
  tick/.style={draw, line width=0.4pt},
  event/.style={anchor=south, align=center, inner sep=1pt},
  span/.style={draw, line width=0.4pt},
  spanlabel/.style={anchor=north, align=center, inner sep=2pt},
  lane/.style={anchor=west, font=\sffamily\scriptsize\itshape},
]

\node[lane] at (0,20)
  {the composer's walk, root field to unreachable field};

% ------------------------------------------------------------ the axis
% Solid while the walk stays inside subgraph sessions; dashed, with the
% arrowhead, for the step it cannot take.
\draw[trunk] (0,0) -- (75,0);
\draw[axis, dashed] (75,0) -- (130,0);

% Ticks 1 and 2: an ordinary stage mark. Tick 3: a cross, not a stroke,
% because the walk does not pass this point.
\foreach \x in {18,75}{\draw[tick] (\x,-1.5) -- (\x,1.5);}
\draw[tick] (128.8,-1.5) -- (131.2,1.5);
\draw[tick] (128.8,1.5) -- (131.2,-1.5);

% ----------------------------------------------------------- the steps
\node[event, text width=30mm] at (18,3)
  {Query.node\\answered by sessions};
\node[event, text width=30mm] at (75,3)
  {... on Session\\still sessions, which\\owns the type};
\node[event, text width=32mm] at (130,3)
  {averageScore\\defined only in ratings};

% ------------------------------------------- the crossing that does not exist
\draw[span] (75,-4) -- (75,-7) -- (130,-7) -- (130,-4);
\node[spanlabel] at (102.5,-8)
  {ratings declares Session with\\resolvable: false, so its key can be\\matched
   against but not entered by:\\no accessible target entity, no route};

\end{tikzpicture}

  \caption{Three steps, and the third does not happen. The first two are
  answered by the same subgraph, which is why the line is unbroken, and the
  third field lives somewhere the composer could find no crossing to. What is
  missing is not a key, since ratings declares one; it is a key the composer is
  allowed to enter through, which is the difference between a seam and a dead
  end and is one argument on one directive.}
  \label{fig:ch07-the-walk}
\end{figure}

\index{resolvability!the reasons are the walk's own record}%
The four reasons underneath the path are the walk's own working, in order: this
is where I started, this is where the field lives, here is why I could not get
from the first to the second, and here is why no other way round exists either.
The third of them is the one that names the fix, and it names it in the
composer's own vocabulary rather than yours: an accessible target entity is a
type with a key you are allowed to enter through.

\subsection{Two Passes, and the Flag That Proves It}

\index{disable-resolvability-validation@\code{--disable-resolvability-validation}!separates the passes}%
The command has an option I would not use and am glad exists:

\begin{minted}[fontsize=\footnotesize,breakanywhere]{text}
wgc router compose -i graph-with-ratings.yaml -o router.json --disable-resolvability-validation
\end{minted}

\index{resolvability!can be switched off, and merging cannot}%
The same two documents, and now they compose. A config is written and a router
would load it. Point the same flag at the shareability collision from
section~\ref{sec:ch07-the-second-document} and it changes nothing at all: that
one still fails, because it was never resolvability's to raise.

\index{composition!is two checks with different powers}%
So composition is two checks, and the flag is the seam between them. Merging is
structural and not optional; a schema whose documents contradict each other is
not a schema. Resolvability is a question about whether a plan exists for every
field, and it can be waived, which is a considerable thing to be able to do by
accident. \code{wgc}'s own help says not to use the flag unless you are
troubleshooting. A graph composed with it is a graph where some field answers at
runtime with an error that no composition ever warned about, and
chapter~\ref{ch:blame-routing} is about how long that takes to trace.

\subsection{Whose Word This Is}

\index{satisfiability!is Apollo's word}%
Apollo calls this satisfiability. Its error reference names a
\code{SATISFIABILITY\_ERROR} and describes it as subgraphs that merge into a
supergraph with \enquote{queries that cannot be satisfied by those
subgraphs}~\autocite{apollofederationerrors}. The Composite Schemas
specification, which chapter~\ref{ch:federation-model} declined to build on,
has a normative section named Validate Satisfiability with an error code of
\code{UNSATISFIABLE\_QUERY\_PATH}, and it writes the check out as an explicit
algorithm over a stack of paths~\autocite{compositeschemasspec}.

\index{resolvability!is Cosmo's word}%
Cosmo uses a different word. Its composition library is its own, with no Apollo
package behind it, and what it calls this check is resolvability: that is the
name on the flag you just read, and it is the name on the source directory and
the error function behind that message. I have not read every page WunderGraph
publishes, so I will not tell you the other word appears nowhere; I will tell
you it appears in none of the places I looked, and that the message on your
screen does not use it.

\index{terminology!the same check under two names}%
Which is worth knowing for the practical reason rather than the pedantic one.
Search for satisfiability after a Cosmo composition has failed and you will
land on Apollo's documentation for a check Apollo's composer runs, described in
terms your error does not use. The idea is the same in both. The words came
from the tool you ran.

\index{resolvability!a walk is my description, not a citation}%
As for calling it a walk: that is my description, and I cite nobody for it
because nobody publishes it in those words. What is published is the algorithm.
The Composite Schemas section above does not describe the check in prose, it
writes it out: start with a stack of one-element paths, take one off, extend
it, push the extensions back, and stop when the stack is
empty~\autocite{compositeschemasspec}. That is a walk whatever anybody calls
it, and it is what the message on the page is a trace of. I would not rest the
claim on the class names in either composer's source, tempting as they are: an
identifier is evidence that somebody once meant something by it and no more
than that.

\index{satisfiability!chapter 15 is the one about reading these}%
One clean failure is all this chapter wants from it.
Chapter~\ref{ch:satisfiability} is the chapter about reading these properly,
including the case that makes them hard: a field reachable four ways, none of
which work, and a composer that blames one of the four.
\index{resolvability!is a second pass|)}%
